%%
% BIThesis 本科毕业设计论文模板 —— 使用 XeLaTeX 编译 The BIThesis Template for Undergraduate Thesis
% This file has no copyright assigned and is placed in the Public Domain.
%%

\begin{appendices}

\newcolumntype{L}[1]{>{\raggedright\arraybackslash}p{#1}}

\section{AI 三级智能提示设计与 Prompt 示例}

从论文整体结构来看，附录真正有必要保留的内容，应当是正文中提到但不适合在主文中过度展开、同时又能直接体现系统设计思路的补充材料。对于本文而言，AI 辅助功能虽然不是论文主体，但它与编码练习场景中的学习支持体验密切相关，且正文已经明确说明该功能被限定在“错误解释与提示”范围内。因此，将 AI 辅助提示设计单独作为附录保留是有必要的，它能够补充说明系统如何在不越过“直接给答案”边界的前提下，为学习者提供分层次的帮助。

需要先说明的是，当前契约与示例材料能够直接证明的事实包括：系统存在 \path{/learner/ai/help} 接口；该接口面向 learner 侧开放；请求中可携带 \texttt{exercise\_id}、\texttt{submission\_id}、\texttt{request\_type}、\texttt{source\_code} 与 \texttt{error\_context\_json}；响应中可返回 \texttt{response\_text}、\texttt{response\_structured\_json}、\texttt{provider\_name}、\texttt{token\_usage}、\texttt{latency\_ms} 和处理状态等字段。契约中已经定义了 \texttt{request\_type} 至少支持 \texttt{error\_explanation} 与 \texttt{hint} 两类请求；mock 数据中还出现了 \texttt{hint\_level = guided} 这一结构化字段。因此，附录中所谓“三级智能提示”，应理解为本文在既有 AI 辅助能力之上的提示分级设计方案，用于展示本系统在学习支持上的可解释设计，而不表述为契约中已经枚举出三种固定 API 类型。

结合编码练习的学习场景，本文将 AI 辅助提示划分为三级：第一级是\textbf{错误定位提示}，目标是帮助学习者理解“哪里出了问题”；第二级是\textbf{修正方向提示}，目标是帮助学习者理解“应该往哪个方向改”；第三级是\textbf{操作建议提示}，目标是帮助学习者在不直接获得完整答案的前提下，得到可以立即尝试的下一步操作。这样的三级结构与移动学习场景相适配：初学者在手机端练习时，最常见的问题并不是完全不会，而是不知道错误产生的位置、原因和第一步该怎么改。

表\ref{tab:ai-three-level-design} 对三级提示的定位进行归纳。

\begin{table}[htbp]
  \centering
  \caption{AI 三级智能提示设计}
  \label{tab:ai-three-level-design}
  \small
  \setlength{\tabcolsep}{4pt}
  \begin{tabular}{L{1.8cm}L{3.0cm}L{3.8cm}L{4.6cm}} 
    \toprule
    提示级别 & 设计目标 & 适用场景 & 输出边界 \\
    \midrule
    一级提示：错误定位 & 帮助学习者识别失败位置与现象 & 刚提交后看到错误、公开用例未通过、布局结果异常 & 说明错误位置、失败现象与相关字段，不直接给出完整修复代码 \\
    二级提示：修正方向 & 帮助学习者理解问题原因与可能的调整方向 & 已知道错误点，但仍不清楚应该修改样式、结构还是断点条件 & 给出原因判断和调整方向，可引用媒体查询、布局属性、选择器覆盖等概念，但不直接替换整段代码 \\
    三级提示：操作建议 & 帮助学习者执行下一步尝试并验证修改效果 & 学习者多次尝试后仍卡住，需要更具体的步骤性帮助 & 给出 2--4 条可执行操作建议，如“先检查断点条件”“再重写 grid-template-columns”，但避免直接输出最终答案 \\
    \bottomrule
  \end{tabular}
\end{table}

在具体实现思路上，三级提示并不是三套完全独立的接口，而更适合作为同一 AI 辅助能力下的三层输出策略。也就是说，移动客户端在调用 \path{/learner/ai/help} 时，仍然可以沿用统一的请求结构；服务端或 AI 编排层根据当前的 \texttt{request\_type}、失败上下文、提交记录和提示次数，决定当前应返回更偏向“解释”、更偏向“方向”，还是更偏向“操作建议”的内容。这样既符合契约优先的接口设计思路，也便于后续在不修改主要接口形状的前提下扩展提示策略。

表\ref{tab:ai-prompt-inputs} 对 AI 提示生成时可利用的输入信息进行归纳。这些输入项均可以从当前契约、示例或页面上下文中获得，不需要依赖额外的隐式状态。

\begin{table}[htbp]
  \centering
  \caption{AI 提示生成的主要输入信息}
  \label{tab:ai-prompt-inputs}
  \small
  \setlength{\tabcolsep}{4pt}
  \begin{tabular}{L{2.6cm}L{3.4cm}L{6.3cm}}
    \toprule
    输入项 & 典型来源 & 作用说明 \\
    \midrule
    题目要求 & 练习详情中的 \texttt{prompt} 字段 & 说明当前练习的目标，例如实现响应式三列卡片布局 \\
    学习者提交代码 & \texttt{source\_code} & 作为 AI 分析错误与给出提示的主要上下文 \\
    提交记录信息 & \texttt{submission\_id}、提交详情与评测状态 & 帮助区分是首次错误还是多次重复失败 \\
    错误上下文 & \texttt{error\_context\_json} & 描述失败用例、错误摘要、断点场景或实际输出 \\
    请求类型 & \texttt{request\_type} & 决定当前是偏向错误解释还是偏向提示引导 \\
    提示层级 & AI 编排规则或结构化返回字段 & 决定当前输出一级、二级还是三级提示风格 \\
    \bottomrule
  \end{tabular}
\end{table}

基于上述输入，本文给出一组适用于论文展示的 Prompt 模板示例。这里的 Prompt 内容用于说明提示设计逻辑，而不是声明论文工程中已将这些模板逐字固化到代码仓库中。附录的目的在于把设计思路展示清楚，使读者能够理解系统为何能够提供“解释而非代做”的 AI 辅助能力。

\subsection*{一级提示 Prompt 模板示例：错误定位}

\begin{quote}
你是一名面向前端初学者的学习辅助教练。请根据题目要求、学习者提交代码和失败上下文，\textbf{只解释当前最可能出错的位置与现象}。\newline
不要直接给出完整答案代码。\newline
请使用简洁中文回答，包含以下三部分：\newline
1) 当前错误最可能出现在哪一段结构或样式；\newline
2) 它导致了什么可观察现象；\newline
3) 学习者下一步应该先检查什么。
\end{quote}

这一层提示适合第一次失败或错误现象较明显的场景。例如，当前示例中的 AI 帮助请求给出了 \texttt{failing\_case = mobile\_single\_column} 和 \texttt{actual\_columns = 2}。在这种情况下，一级提示不必解释完整的 CSS 布局原理，只需指出“窄屏断点下仍保留了两列布局，优先检查媒体查询是否生效，以及对应容器是否仍在使用原有列定义”即可。

\subsection*{二级提示 Prompt 模板示例：修正方向}

\begin{quote}
你是一名面向前端初学者的学习辅助教练。请在不直接给出完整代码答案的前提下，\textbf{解释问题产生的原因，并给出修正方向}。\newline
请结合题目要求、提交代码和失败上下文，输出以下内容：\newline
1) 失败最可能由哪一类原因导致；\newline
2) 学习者应优先修改布局结构、选择器覆盖、断点条件还是属性值；\newline
3) 给出 2 条不含完整答案的方向性建议。
\end{quote}

二级提示比一级提示更进一步，它不只告诉学习者“哪里错了”，还告诉学习者“为什么会这样”。例如在响应式卡片布局练习中，如果移动端仍然显示两列，那么二级提示可以指出：问题通常来自媒体查询断点未命中、媒体查询中未覆盖原有布局属性，或者修改了错误的选择器。这样的输出能够帮助学习者从单次报错中形成更稳定的前端调试思路。

\subsection*{三级提示 Prompt 模板示例：操作建议}

\begin{quote}
你是一名面向前端初学者的学习辅助教练。学习者已经获得基础解释，但仍然没有修复问题。请输出\textbf{可执行的下一步操作建议}，帮助其继续尝试。\newline
要求：\newline
1) 仅给出步骤，不直接生成完整答案；\newline
2) 每一步都应能在当前编辑器中立即尝试；\newline
3) 最多给出 4 步；\newline
4) 最后提醒学习者重新提交并观察哪个公开用例发生变化。
\end{quote}

三级提示适合学习者在多次尝试后仍然卡住的情况。它最接近“教学支架”，但仍然不能越过论文中设定的边界。以当前示例为例，三级提示可以设计成：先检查媒体查询断点是否与测试场景一致；再确认媒体查询内是否重写了容器的列定义；接着在本地预览中缩小视口观察列数变化；最后重新提交并关注公开测试是否通过。这样的设计既能提高学习者继续尝试的成功率，也能避免 AI 直接替代学习者完成编码任务。

综合来看，保留 AI 提示设计附录是有必要的，原因不在于该功能本身是否庞大，而在于它能够直接回应论文中的一个关键设计判断：系统的 AI 辅助能力不是面向“自动写代码”，而是面向“在练习受阻时提供分层提示”。这一点若只在正文中概括描述，容易显得抽象；通过附录列出提示层级与 Prompt 模板，能够使该设计更加具体、可解释，也更符合毕业论文附录“补充主文论述而不打断正文节奏”的用途。

\section{OpenAPI 接口列表与对应功能表单}

除了 AI 提示设计外，另一类有保留价值的附录是“OpenAPI 接口列表与对应功能表单”。这类附录的必要性在于：论文正文已经多次强调系统采用契约优先设计，且移动客户端、服务端与管理后台端围绕同一份 OpenAPI 文档协作；如果附录能够进一步把“接口”和“页面/表单”一一对应起来，就能更直接说明契约并非抽象说明书，而是具体指导了登录、资料维护、练习提交、AI 辅助以及后台内容维护等页面功能的字段设计。

在本论文范围内，没有必要把全部 OpenAPI 路径都列为附录。更合适的做法，是只选择那些既与正文高频功能相关、又能代表不同页面表单类型的接口：例如登录表单、个人资料表单、练习提交表单、AI 辅助请求表单、课程创建表单、挑战创建表单和练习创建表单。这样既能体现契约优先的工程方法，也不会把附录写成冗长的接口手册。

表\ref{tab:appendix-openapi-forms} 对关键接口与对应表单进行归纳。

\begin{table}[htbp]
  \centering
  \caption{关键 OpenAPI 接口与对应功能表单}
  \label{tab:appendix-openapi-forms}
  \small
  \setlength{\tabcolsep}{4pt}
  \begin{tabular}{L{3.0cm}L{1.0cm}L{3.3cm}L{5.0cm}}
    \toprule
    接口路径 & 方法 & 对应页面或表单 & 主要表单字段 \\
    \midrule
    \path{/auth/learner/login} & POST & 学习者登录表单 & email、password、device\_id、device\_name、platform \\
    \path{/learner/profile} & PATCH & 学习者个人资料与偏好表单 & nickname、avatar\_url、bio、theme\_mode、daily\_goal\_minutes \\
    \path{/learner/submissions} & POST & 练习提交表单 & exercise\_id、source\_code、content\_version、rule\_version \\
    \path{/learner/ai/help} & POST & AI 辅助请求表单 & exercise\_id、submission\_id、request\_type、source\_code、error\_context\_json \\
    \path{/admin/courses} & POST & 管理后台端课程创建表单 & course\_code、title、summary、description、cover\_image\_url、difficulty、estimated\_minutes、status、sort\_order、content\_version、published\_at \\
    \path{/admin/challenges} & POST & 管理后台端挑战创建表单 & challenge\_code、title、summary、related\_course\_id、difficulty、reward\_xp、status、sort\_order、content\_version、stages \\
    \path{/admin/exercises} & POST & 管理后台端练习创建表单 & chapter\_id、exercise\_code、title、prompt、exercise\_type、starter\_code、language、difficulty、pass\_score、max\_attempts\_per\_day、status、content\_version、options、test\_cases \\
    \bottomrule
  \end{tabular}
\end{table}

为了进一步说明这些接口如何映射到真实页面表单，表\ref{tab:appendix-form-purpose} 从“功能意图”角度进行补充。

\begin{table}[htbp]
  \centering
  \caption{功能表单的作用与表单意图说明}
  \label{tab:appendix-form-purpose}
  \small
  \setlength{\tabcolsep}{4pt}
  \begin{tabular}{L{2.7cm}L{4.0cm}L{5.7cm}}
    \toprule
    表单名称 & 主要作用 & 设计说明 \\
    \midrule
    学习者登录表单 & 建立认证会话并绑定移动端设备上下文 & 除账号密码外，还包含设备标识、设备名称与平台信息，说明登录不仅是身份校验，还服务于会话管理与设备侧上下文识别 \\
    学习者个人资料表单 & 维护昵称、头像、简介与学习偏好 & 字段较少但直接影响个人中心展示和个性化偏好，是移动客户端中典型的轻量编辑表单 \\
    练习提交表单 & 将当前代码与版本信息提交到评测流程 & 不只是提交源码，还要附带内容版本与规则版本，以保证评测结果与当前练习配置一致 \\
    AI 辅助请求表单 & 把当前卡点上下文发送给 AI 辅助能力 & 同时包含题目、提交和错误上下文信息，体现 AI 帮助不是孤立问答，而是依附于当前练习状态的上下文化提示请求 \\
    课程创建表单 & 由管理后台端创建课程基础信息 & 既包含标题、摘要、封面等展示字段，也包含业务编码、发布状态和版本号等后台维护字段 \\
    挑战创建表单 & 由管理后台端配置挑战与关卡结构 & 除基础信息外，还包含关卡数组、奖励值和解锁规则，体现挑战表单比普通内容表单更强调状态链与阶段结构 \\
    练习创建表单 & 由管理后台端配置练习题干、选项和测试用例 & 表单最复杂，既要配置题目本体，也要配置公开/隐藏测试用例与通过阈值，是契约优先设计最能体现结构价值的典型场景 \\
    \bottomrule
  \end{tabular}
\end{table}

从这些表单可以看出，OpenAPI 对系统的作用不仅是给出一组接口路径，更是提前定义了页面表单应当如何组织字段。以学习者登录表单为例，示例中除 \texttt{email} 和 \texttt{password} 外，还要求提交 \texttt{device\_id}、\texttt{device\_name} 和 \texttt{platform}。这说明移动客户端登录页在设计时就需要考虑设备上下文，而不是只做最简的账号密码输入。又如个人资料更新接口中出现 \texttt{theme\_mode} 与 \texttt{daily\_goal\_minutes} 字段，意味着个人资料页并不只是静态信息展示，还承担了偏好设置入口的作用。

在练习与 AI 辅助场景中，契约与表单之间的对应关系更为明显。练习提交表单要求提交 \texttt{exercise\_id}、\texttt{source\_code}、\texttt{content\_version} 和 \texttt{rule\_version}，这意味着练习页面不仅要保存当前代码内容，还要与当前题目版本和评测规则版本保持一致。AI 辅助请求表单则进一步扩展了上下文维度，除练习标识与提交标识外，还需要提供 \texttt{request\_type} 和 \texttt{error\_context\_json}。这说明 AI 帮助入口并不是一个脱离练习状态的独立聊天框，而是紧密依附于“当前哪道题、哪次提交、当前报什么错”的上下文化辅助入口。

在管理后台端，课程、挑战和练习创建表单则体现了契约优先设计的另一层意义，即通过结构化字段约束后台内容维护行为。课程创建表单需要填写标题、摘要、描述、封面、难度、状态和发布时间等基础字段；挑战创建表单进一步引入 \texttt{stages}、星级规则和解锁规则；练习创建表单则最为复杂，既要配置题目描述与起始代码，又要配置选项和测试用例。若没有 OpenAPI 与示例文件约束，这些复杂表单极易在字段命名、数据层级或必填规则上出现前后不一致；而附录把接口与表单对应后，读者可以更直观看出契约如何实质性地参与了后台页面设计。

因此，从论文附录的必要性来看，保留“OpenAPI 接口列表与对应功能表单”是合理的，因为它能够把正文中关于契约优先设计的论述落到具体页面和字段层面。相反，若附录继续保留大而泛的跨端模块映射、测试矩阵和联调场景，则容易与正文已有内容重复，且不如这类“接口--表单”对应关系更直接、更有辨识度。按照当前论文的重点，附录保留 AI 提示设计和 OpenAPI/表单对应这两类内容，已经足以形成清晰、必要且有区分度的补充材料。

\end{appendices}
